\section{Лемма Дюбуа-Раймона. Уравнения Эйлера-Лагранжа для простейшей
вариационной задачи с закрепленными концами}
\textbf{Введение}\\
Пусть $I[u]$ --- \textit{функционал энергии}, который сопоставляет каждой 
допустимой функции $u(x)$ из некоторого пространства подмножества $\R^n$
вещественное число. \\
Для нахождения функции $u(x)$, доставляющей минимум функционалу $I[u]$,
применяется подход, аналогичный поиску экстремумов функций конечного числа
переменных в математическом анализе. Роль обычной производной здесь играет
\textit{первая вариация} функционала,
обозначаемая как $I'[u]$.\\
Введем оператор первой вариации $A[u]$, положив по определению:
\begin{equation}
    A[u] = I'[u].
\end{equation}
Тогда соотношение
\begin{equation}
    A[u] = 0 
\end{equation}
задает \textit{необходимое условие экстремума} функционала $I[u]$.\\

\textbf{Билет}\\
Рассмотрим простейшую вариационную задачу для функционала энергии $I[u]$. 
Требуется найти минимум функционала:
\begin{equation}
    I[w] = \int_U L(Dw(x), w(x), x) \, dx \rightarrow \min,
\end{equation}
где $U \subset \R^n$ --- ограниченное открытое множество с гладкой границей
$\partial U$. \\
Лагранжиан представляет собой гладкую функцию 
$$L \colon \R^n \times \R \times \overline{U} \to \R, \quad
L(p, z, x) = L(p_1, \dots, p_n, z, x_1, \dots, x_n)$$
Будем искать минимум в классе гладких функций
с закрепленными концами (с заданными непрерывными значениями на границе):
\begin{equation}
    \mathcal{M} = \{ w \in C^\infty(U) \mid w|_{\partial U} = y \in C(\overline{U}) \}.
\end{equation}
% \begin{definition}
%     Введем векторные обозначения для производных лагранжиана по соответствующим 
%     переменным:
%     \begin{equation}
%         D_p L = (L_{p_1}, \dots, L_{p_n}), 
%         \quad D_z L = L_z, 
%         \quad D_x L = (L_{x_1}, \dots, L_{x_n}).
%     \end{equation}
% \end{definition}

\begin{lemma}[Дюбуа-Раймона]
    Пусть функция $f(x)$ непрерывна. 
    Если для любой гладкой функции с комрактным носителем
    $v(x) \in C_0^\infty(U)$ выполняется интегральное равенство
    \begin{equation}
        \int_U f(x) v(x) \, dx = 0,
    \end{equation}
    то подинтегральная функция тождественно равна нулю: 
    $f(x) \equiv 0$ для всех $x \in U$.
\end{lemma}

\begin{theorem}[Уравнение Эйлера-Лагранжа]
    Если функция $w(x) \in \mathcal{M}$ доставляет минимум функционалу $I[w]$,
    то она удовлетворяет уравнению Эйлера-Лагранжа 
    (где $A[w] = I'[w] = 0$ --- условие равенства нулю первой вариации):
    \begin{equation}
        - \sum_{i=1}^n \frac{\partial}{\partial x_i} L_{p_i} (Dw, w, x) + L_z(Dw, w, x) = 0.
    \end{equation}
\end{theorem}

\begin{proof}
    Пусть $w(x)$ --- минимизирующая функция. Фиксируем произвольную функцию $v(x) \in C_0^\infty(U)$ (гладкую функцию с компактным носителем)[cite: 21, 30, 188].
    Рассмотрим вспомогательную функцию одной переменной,
    задающую значение функционала на однопараметрическом семействе:
    \begin{equation}
        i(\tau) = I[w + \tau v].
    \end{equation}
    Поскольку $w(x)$ доставляет минимум, необходимое условие экстремума
    имеет вид $\frac{\partial i(\tau)}{\partial \tau} \big|_{\tau = 0} = 0$.
    Вычислим эту производную:
    \begin{equation}
        \frac{\partial i(\tau)}{\partial \tau} = \int_U \left\{ \sum_{i=1}^n L_{p_i}(Dw + \tau Dv, w + \tau v, x) v_{x_i} + L_z(Dw + \tau Dv, w + \tau v, x) v \right\} dx.
    \end{equation}
    Полагая $\tau = 0$ и приравнивая выражение к нулю, 
    получаем первую вариацию функционала:
    \begin{equation}
        \int_U \left\{ \sum_{i=1}^n L_{p_i}(Dw, w, x) v_{x_i}(x) + L_z(Dw, w, x) v(x) \right\} dx = 0.
    \end{equation}
    Проинтегрируем первое слагаемое по частям. 
    Так как функция $v(x)$ с компактным носителем, 
    внеинтегральные граничные члены обращаются в ноль, получим:
    \begin{equation}
        \int_U \left\{ - \sum_{i=1}^n \frac{\partial}{\partial x_i} L_{p_i}(Dw, w, x) + L_z(Dw, w, x) \right\} v(x) \, dx = 0.
    \end{equation}
    Поскольку это равенство выполнено для любой пробной функции 
    $\forall v(x) \in C_0^\infty(U)$, в силу леммы Дюбуа-Раймона выражение в
    фигурных скобках должно быть тождественно равно нулю.
    Отсюда следует дифференциальное уравнение Эйлера-Лагранжа:
    \begin{equation}
        - \sum_{i=1}^n \frac{\partial}{\partial x_i} L_{p_i} (Dw, w, x) + L_z(Dw, w, x) = 0.
    \end{equation}
\end{proof}

\subsection{Пример 1. Принцип Дирихле}
Рассмотрим простейший лагранжиан, зависящий только от квадрата нормы градиента 
функции: $L(p,z,x) = \frac{1}{2}(p,p) = \frac{1}{2}|p|^2$. В этом случае 
частные производные лагранжиана принимают вид $L_{p_i} = p_i$ и $L_z = 0$.

Соответствующий функционал энергии (называемый интегралом Дирихле) записывается как:
\begin{equation}
    I[w] = \frac{1}{2} \int_U |Dw|^2 \, dx \rightarrow \min
\end{equation}

Подставляя производные в общую схему, получаем, что уравнение Эйлера-Лагранжа 
переходит в классическое уравнение Лапласа:
\begin{equation}
    -\Delta w = 0
\end{equation}

\subsection{Пример 2. Обобщенный принцип Дирихле}
Пусть лагранжиан задан произвольной симметричной квадратичной формой по 
переменным градиента со свободным линейным членом:
\begin{equation}
    L(p,z,x) = \frac{1}{2}\sum_{i,j=1}^n a^{ij}(x)p_i p_j - z f(x)
\end{equation}
где матрица коэффициентов $a^{ij}(x)$ предполагается симметричной 
($a^{ij} = a^{ji}$).

Функционал задачи имеет вид:
\begin{equation}
    I[w] = \int_U \left\{ \frac{1}{2}\sum_{i,j=1}^n a^{ij}(x)w_{x_i}(x)w_{x_j}(x) - w(x)f(x) \right\} dx \rightarrow \min
\end{equation}

Находя частные производные лагранжиана, имеем 
$L_{p_i} = \sum_{j=1}^n a^{ij}(x)p_j$ и $L_z = -f(x)$. 
Уравнение Эйлера-Лагранжа в данном случае представляет собой линейное
дифференциальное уравнение второго порядка в дивергентной форме 
(уравнение Пуассона):
\begin{equation}
    -\sum_{i,j=1}^n \frac{\partial}{\partial x_i} \left( a^{ij}(x) w_{x_j} \right) = f(x)
\end{equation}

\subsection{Пример 3. Уравнение Пуассона}
Пусть $f \colon \R \rightarrow \R$ --- заданная гладкая функция, и определена 
её первообразная $F(z) = \int_0^z f(s)\,ds$. Рассмотрим лагранжиан:
\begin{equation}
    L(p,z,x) = \frac{1}{2}|p|^2 - F(z)
\end{equation}

Функционал энергии записывается в виде:
\begin{equation}
    I[w] = \int_U \left\{ \frac{1}{2}|Dw|^2 - F(w(x)) \right\} dx \rightarrow \min
\end{equation}

Производные лагранжиана равны $L_{p_i} = p_i$ и $L_z = -F'(z) = -f(z)$. 
Подстановка этих выражений в уравнение Эйлера-Лагранжа приводит к:
\begin{equation}
    \Delta w(x) = f(w(x))
\end{equation}

\subsection{Пример 4. Минимальные поверхности}
В качестве еще одного классического примера в рамках простейшей вариационной 
задачи выступают:

1. \textbf{Задача о минимальных поверхностях}. \\
Для нахождения графика функции $w(x)$, минимизирующего площадь поверхности 
при заданном контуре, лагранжиан имеет вид:
\begin{equation}
    L(p,z,x) = \sqrt{1 + |p|^2}
\end{equation}
Тогда $L_{p_i} = \frac{p_i}{\sqrt{1+|p|^2}}$, 
а уравнение Эйлера-Лагранжа задает среднюю кривизну поверхности:
\begin{equation}
    \sum_{i=1}^n \frac{\partial}{\partial x_i} \left( \frac{w_{x_i}}{\sqrt{1 + |Dw|^2}} \right) = 0
\end{equation}

2. \textbf{Пример Гильберта}. \\
Одномерная задача, иллюстрирующая особенности вариационных методов, 
где функционал имеет вид:
\begin{equation}
    I[x(t)] = \int_0^1 t^{2/3} \left( \frac{dx}{dt} \right)^2 dt 
    \rightarrow \min, \quad x(0)=0, \, x(1)=1
\end{equation}
Соответствующее уравнение Эйлера-Лагранжа 
$\frac{d}{dt}\left( t^{2/3} \frac{dx}{dt} \right) = 0$ 
имеет экстремаль $\hat{x}(t) = t^{1/3}$ 
(но при этом экстремаль не является непрерывно дифференцируемой).
